    \documentclass[11pt]{article}
    
    % --- UNIVERSAL PREAMBLE BLOCK ---
    % Set 4:5 Aspect Ratio (Standard LinkedIn Portrait Carousel 1080x1350 equivalent)
    \usepackage[paperwidth=8in, paperheight=10in, top=1.2in, bottom=1.2in, left=1in, right=1in]{geometry}
    \usepackage{fontspec}
    \usepackage[english, bidi=basic, provide=*]{babel}
    
    \babelprovide[import, onchar=ids fonts]{english}
    
    % Set premium typography using Noto Sans system fonts
    \babelfont{rm}{Noto Sans}
    \babelfont{sf}{Noto Sans}
    \babelfont{tt}{Noto Sans Mono}
    
    % --- PACKAGES FOR MODERN DESIGN & MATHS ---
    \usepackage{amsmath}
    \usepackage{amssymb}
    \usepackage{xcolor}
    \usepackage{graphicx}
    \usepackage{listings}
    \usepackage{mdframed}
    \usepackage{tikz}
    \usetikzlibrary{shapes.geometric, arrows.meta, positioning, shadows.blur}
    \usepackage{tcolorbox}
    \tcbuselibrary{listings, skins, breakable}
    \usepackage{fancyhdr}
    \usepackage{enumitem}
    
    % --- DESIGN TOKENS (Stripe & Linear Inspired) ---
    \definecolor{brandnavy}{HTML}{0A2540}    % Deep Primary Navy
    \definecolor{brandgray}{HTML}{425466}    % Slate Gray
    \definecolor{brandlight}{HTML}{F8F9FA}   % Light Soft Gray
    \definecolor{brandblue}{HTML}{635BFF}    % Accent Classic Blue
    \definecolor{brandgreen}{HTML}{00D4B2}   % Success/Muted Green
    \definecolor{brandorange}{HTML}{F25C05}  % Accent Orange
    \definecolor{brandred}{HTML}{E5484D}     % Warning Red
    
    % Premium Dark Theme for Code Blocks
    \definecolor{codebackground}{HTML}{141419}
    \definecolor{codekeyword}{HTML}{FF7B72}
    \definecolor{codestring}{HTML}{A5D6FF}
    \definecolor{codecomment}{HTML}{8B949E}
    \definecolor{codetext}{HTML}{E6EDF3}
    \definecolor{codeframe}{HTML}{30363D}
    
    % --- ZERO PARAGRAPH INDENT ---
    \setlength{\parindent}{0pt}
    \setlength{\parskip}{1.2em}
    
    % Custom TColorBox for styled content blocks
    \newtcolorbox{infobox}[1][]{
        colback=brandlight, colframe=brandblue, arc=5pt, boxrule=0.8pt,
        left=15pt, right=15pt, top=12pt, bottom=12pt, #1
    }
    
    \newtcolorbox{darkbox}[1][]{
        colback=brandnavy, colframe=brandnavy, arc=5pt, boxrule=0.8pt,
        left=15pt, right=15pt, top=12pt, bottom=12pt, 
        coltext=white, #1
    }
    
    % --- HEADER & FOOTER CONFIGURATION ---
    \pagestyle{fancy}
    \fancyhf{}
    \renewcommand{\headrulewidth}{0pt}
    \renewcommand{\footrulewidth}{0.5pt}
    
    \fancyhead[L]{\fontsize{9}{11}\selectfont\color{brandgray}\MakeUppercase{First Principles Deep Dive}}
    \fancyhead[R]{\fontsize{9}{11}\selectfont\color{brandblue}\textbf{AI ARCHITECTURE}}
    
    \fancyfoot[L]{\fontsize{8}{10}\selectfont\color{brandgray}
    Author: Dibayendu Mukherjee \hspace{0.7em}$\cdot$\hspace{0.7em}
    \mbox{AI Research Engineer \hspace{0.7em}\textbullet\hspace{0.7em} Systems Engineer}}
    
    \fancyfoot[R]{\fontsize{8}{10}\selectfont\color{brandgray}Slide \thepage}
    
    \begin{document}
    
    % ==========================================================
    % SLIDE 1: COVER
    % ==========================================================
    \thispagestyle{empty}
    \vspace*{0.8in}
    
    \begin{center}
    
        {\fontsize{14}{16}\selectfont\color{brandblue}\textbf{
        THE MECHANICS OF MEMORY $\cdot$ \quad SEQUENCE MODELS
        }}
    
        \vspace{0.8in}
    
        {\fontsize{40}{48}\selectfont\color{brandnavy}\bfseries
        How does AI\\
        know what\\
        to remember?
        }
    
        \vspace{0.6in}
    
        {\fontsize{18}{24}\selectfont\color{brandgray}
        Understanding memory flow, gates, and the\\
        architecture behind the Gated Recurrent Unit (GRU).
        }
    
        \vfill
    
        \begin{minipage}{0.8\textwidth}
            \centering
            \rule{0.6\linewidth}{0.5pt}
    
            \vspace{0.2in}
    
            {\fontsize{12}{14}\selectfont\color{brandnavy}
            \textbf{Dibayendu Mukherjee}
            }
    
            \vspace{0.05in}
    
            {\fontsize{10}{12}\selectfont\color{brandgray}
            AI Research Engineer \quad $\cdot$ \quad Systems Engineer
            }
    
        \end{minipage}
    
    \end{center}
    
    \newpage
    
    % ==========================================================
    % SLIDE 2: THE PROBLEM
    % ==========================================================
    \vspace*{0.4in}
    
    \begin{center}
    {\fontsize{26}{30}\selectfont\color{brandnavy}\bfseries 
    The Problem with\\Traditional RNNs}
    \end{center}
    
    \vspace{0.3in}
    
    {\large\color{brandgray}
    In my last post, we explored the shift from LSTMs to Transformers.
    }
    
    \vspace{0.15in}
    
    But while building and studying sequence models, I kept returning to a
    fundamental question:
    
    \textbf{Why do simple RNNs struggle to remember the past?}
    
    \vspace{0.15in}
    
    Traditional RNNs face a major limitation:
    
    \textbf{long-term dependencies.}
    
    \vspace{0.1in}
    
    With every new token, the hidden state is repeatedly updated. 
    Important information from earlier steps slowly fades as new information arrives.
    
    \vspace{0.15in}
    
    \begin{darkbox}
    
    {\large\textbf{The Realization}}
    
    \vspace{0.05in}
    
    As sequences become longer, an RNN has no way to decide:
    
    \begin{center}
    \textit{"Should I keep this information, or replace it?"}
    \end{center}
    
    Without a mechanism to control memory flow, useful signals get diluted by noise.
    
    \end{darkbox}
    
    \vfill
    \newpage
    
    % ==========================================================
    % SLIDE 3: THE INTUITION
    % ==========================================================
    \section*{The Intuition: Reading an Article}
    
    \vspace{0.15in}
    
    Imagine reading a long article and trying to summarize it.
    
    Your brain does not store every word. It creates a compressed memory of what matters.
    
    \vspace{0.2in}
    
    \begin{minipage}[t]{0.45\textwidth}
        {\large\color{brandnavy}\textbf{Keep in Memory}}
    
        \begin{itemize}[leftmargin=*, itemsep=0.1in, label={\color{brandblue}\textbullet}]
            \item The main idea.
            \item Important events.
            \item Meaningful relationships.
        \end{itemize}
    
    \end{minipage}
    \hfill
    \begin{minipage}[t]{0.45\textwidth}
    
        {\large\color{brandgray}\textbf{Discard}}
    
        \begin{itemize}[leftmargin=*, itemsep=0.1in, label={\color{brandorange}\textbullet}]
            \item Repeated information.
            \item Irrelevant details.
            \item Temporary noise.
        \end{itemize}
    
    \end{minipage}
    
    \vspace{0.3in}
    
    \begin{infobox}
    
    A sequence model needs the same ability.
    
    At every step, it must decide:
    
    \vspace{0.1in}
    
    \begin{center}
    \textbf{"What should I forget?"} \quad
    \textbf{"What should I keep?"} \quad
    \textbf{"What new information should I add?"}
    \end{center}
    
    \end{infobox}
    
    \vfill
    \newpage
    
    % ==========================================================
    % SLIDE 4: THE SOLUTION (GRU)
    % ==========================================================
    
    \section*{The Solution: Controlled Memory}
    
    \vspace{0.15in}
    
    This is exactly why gates were introduced.
    
    Enter the \textbf{Gated Recurrent Unit (GRU)}.
    
    \vspace{0.01in}
    
    A GRU does not force the network to remember everything.
    
    Instead, it learns \textbf{when information should flow and when it should stop.}
    
    \vspace{0.01in}
    
    These learned gates continuously answer:
    
    \begin{center}
    \textit{"Should this information update my memory?"}
    \end{center}
    
    \vspace{0.01in}
    
    \begin{center}
        \includegraphics[width=0.85\textwidth]{gru_architecture.png}
    \end{center}
    
    \vspace{0.2in}
    
    Behind this simple diagram are mathematical gates that regulate information flow, allowing the network to preserve important patterns across long sequences.
    
    \vfill
    \newpage
    
% ==========================================================
% SLIDE 5: INSIDE THE GATES
% ==========================================================
\section*{Inside the Gates:}

\vspace{0.05in}

\textbf{The Evolution of Memory Control:}
\begin{itemize}[leftmargin=*, itemsep=0.02in, label={\color{brandblue}\textbullet}]
    \item \textbf{RNN Problem:} Blindly overwrites memory with every new input.
    \item \textbf{LSTM Solution:} Introduced gates to protect memory, but added high complexity (two memory states, three gates).
    \item \textbf{GRU Simplification:} \textit{"Can we keep the same ability but simplify the mechanism?"}
\end{itemize}

\vspace{0.05in}

Unlike an LSTM which splits memory into a cell state and hidden state, GRU elegantly merges them into a single state ($h_t$), controlled by two core decisions:

\vspace{0.05in}

\begin{minipage}[t]{0.45\textwidth}
    {\large\color{brandnavy}\textbf{Update Gate ($z_t$)}}

    \textit{"Should I keep the past or replace it?"}

    \vspace{0.02in}
    Controls the blending of old memory and new information.
\end{minipage}
\hfill
\begin{minipage}[t]{0.45\textwidth}
    {\large\color{brandnavy}\textbf{Reset Gate ($r_t$)}}

    \textit{"When creating new memory, should I look at history?"}

    \vspace{0.02in}
    Controls how much past context influences the candidate memory.
\end{minipage}

\vspace{0.05in}

\begin{center}
    \includegraphics[width=0.55\textwidth]{gru_equations.png}
\end{center}

\vspace{0.05in}

\begin{darkbox}
{\large\textbf{The Mathematical Shift}}

The GRU mathematically simplifies the architecture by combining the roles of an LSTM's forget and input gates into a single update gate ($z_t$). The network \textit{learns} this mixing ratio through gradient descent.
\end{darkbox}

% ==========================================================
% SLIDE 6: THE EQUATIONS EXPLAINED
% ==========================================================
\section*{Breaking Down the Math}

\vspace{0.05in}

{\fontsize{10}{12}\selectfont\color{brandgray}
The equations are not just formulas, they are three learned mechanisms that control how information flows through memory.
}

\vspace{0.15in}

% --- UPDATE GATE ---
\begin{darkbox}[top=8pt, bottom=8pt, left=12pt, right=12pt]
{\large\color{brandlight}\textbf{1. Update Gate ($z_t$)}} \quad $z_t=\sigma(W_z[h_{t-1},x_t]+b_z)$

\vspace{0.05in}
{\fontsize{10}{12}\selectfont
\textit{"How much of my previous memory should I keep?"} \\
The GRU creates one vector containing \textbf{old memory + current information}: $[h_{t-1},x_t]$. This passes through a learned transformation and a sigmoid function ($0 \leq z_t \leq 1$), acting as a soft memory controller:
\begin{itemize}[leftmargin=*, itemsep=0pt, parsep=0pt, topsep=3pt, label={\color{brandlight}\textbullet}]
    \item $z_t \approx 1 \rightarrow$ \textbf{preserve old memory}
    \item $z_t \approx 0 \rightarrow$ \textbf{replace with new information}
\end{itemize}
}
\end{darkbox}

\vspace{0.1in}

% --- RESET GATE ---
\begin{darkbox}[top=8pt, bottom=8pt, left=12pt, right=12pt]
{\large\color{brandlight}\textbf{2. Reset Gate ($r_t$)}} \quad $r_t=\sigma(W_r[h_{t-1},x_t]+b_r)$

\vspace{0.05in}
{\fontsize{10}{12}\selectfont
\textit{"When creating new memory, how much past context should I use?"} \\
It decides how much previous information should influence the candidate memory by controlling the previous hidden state: $r_t\odot h_{t-1}$. This allows the GRU to forget irrelevant history when forming new memories:
\begin{itemize}[leftmargin=*, itemsep=0pt, parsep=0pt, topsep=3pt, label={\color{brandlight}\textbullet}]
    \item $r_t \approx 0 \rightarrow$ \textbf{ignore previous context}
    \item $r_t \approx 1 \rightarrow$ \textbf{use previous context}
\end{itemize}
}
\end{darkbox}

\vspace{0.1in}

% --- CANDIDATE & BLEND ---
\begin{infobox}[top=8pt, bottom=8pt, left=12pt, right=12pt]
{\large\color{brandnavy}\textbf{3. Candidate Memory + Final Hidden State}}

\vspace{0.05in}
{\fontsize{10}{12}\selectfont
\begin{minipage}[t]{0.48\textwidth}
\textbf{Candidate Memory ($g_t$):} \\
$g_t=\tanh(W_g[r_t\odot h_{t-1},x_t]+b_g)$ \\
\textit{"What new information should I remember?"}

\vspace{0.05in}
The reset gate filters old memory first. Then the GRU creates a candidate memory by combining \textbf{filtered history + current input}.
\end{minipage}
\hfill
\begin{minipage}[t]{0.48\textwidth}
\textbf{Final Blend ($h_t$):} \\
$h_t=z_t\odot h_{t-1} + (1-z_t)\odot g_t$ \\
\textit{Linear Interpolation}

\vspace{0.05in}
The $\odot$ denotes element-wise multiplication. Here, \textbf{$z_t$} is the percentage of old memory to keep, and \textbf{$(1-z_t)$} is the exact remaining percentage used for the new candidate. This perfectly balances past and present, preventing memory from exploding.
\end{minipage}
}
\end{infobox}

\vfill
\newpage

% ==========================================================
% SLIDE 7: GRU vs LSTM
% ==========================================================
\section*{GRU vs LSTM:\\Choosing the Right Memory}

\vspace{0.05in}

Both GRU and LSTM are designed to preserve long‑term
dependencies and mitigate the vanishing gradient problem.
The key difference lies in how they represent and manage memory.

\vspace{0.15in}

% --- Equally distributed table with booktabs ---
{
\small
\renewcommand{\arraystretch}{1.2}
\begin{center}
\begin{tabular}{@{}p{\dimexpr(\textwidth-24pt)/3\relax}
                p{\dimexpr(\textwidth-24pt)/3\relax}
                p{\dimexpr(\textwidth-24pt)/3\relax}@{}}
\toprule
\rowcolor{brandlight}
\textbf{Feature} &
\color{brandblue}\textbf{LSTM} &
\color{brandorange}\textbf{GRU} \\
\midrule

\textbf{Memory Representation} &
\textbf{Cell State ($C_t$):} Long-term memory

\textbf{Hidden State ($h_t$):} Current working memory &
\textbf{Hidden State ($h_t$):}
Stores both long‑term memory and current context \\

\textbf{Memory Structure} &
Two separate memory streams &
One unified memory stream \\

\textbf{Memory Control} &
More explicit and flexible &
Simplified and efficient \\

\textbf{Gates} &
Forget, Input, Output &
Update, Reset \\

\textbf{Parameters} &
More parameters &
Fewer parameters \\

\textbf{Training Speed} &
Usually slower &
Usually faster \\

\textbf{Data Efficiency} &
Often benefits from larger datasets &
Often performs well with smaller datasets \\

\textbf{Performance} &
Can perform better on complex sequential tasks &
Often comparable with lower computational cost \\

\bottomrule
\end{tabular}
\end{center}
}

\vspace{0.1in}

% --- Explanation box ---
\begin{infobox}[top=8pt, bottom=8pt, left=12pt, right=12pt]
{\large\color{brandnavy}\textbf{Think of it this way:}}

\vspace{0.05in}
{\fontsize{10}{12}\selectfont
\begin{itemize}[leftmargin=*, itemsep=0pt, parsep=0pt, topsep=3pt, label={\color{brandnavy}\textbullet}]
  \item \textbf{LSTM} keeps two memories:
    \begin{itemize}[leftmargin=*, itemsep=0pt, parsep=0pt, topsep=3pt, label={\color{brandnavy}\textbullet}]
      \item Cell State ($C_t$): A long‑term notebook storing important information across many time steps.
      \item Hidden State ($h_t$): A short‑term workspace used to produce the current output.
    \end{itemize}
  \item \textbf{GRU} merges both roles into a single Hidden State ($h_t$), making the architecture simpler, faster, and more computationally efficient.
\end{itemize}
}
\end{infobox}
\newpage

% ==========================================================
% SLIDE 8: DECISION MAKING
% ==========================================================

\section*{When to Choose Which?}

\vspace{0.1in}

There is no universally better architecture.
The right choice depends on the complexity of the task,
available data, and computational constraints.

\vspace{0.15in}

\begin{minipage}[t]{0.47\textwidth}

{\large\color{brandorange}\textbf{Choose GRU}}

\vspace{0.05in}

\begin{itemize}[leftmargin=*, itemsep=0.08in,
label={\color{brandorange}\textbullet}]

\item You want a simpler architecture.

\item Training speed and inference latency are important.

\item Memory or computational resources are limited.

\item The dataset is relatively small.

\item Performance is comparable to LSTM for your task.

\end{itemize}

\vspace{0.08in}

\textbf{Common Applications}

{\small
IoT forecasting\\
Speech recognition\\
Edge AI\\
Mobile applications
}

\end{minipage}
\hfill
\begin{minipage}[t]{0.47\textwidth}

{\large\color{brandblue}\textbf{Choose LSTM}}

\vspace{0.05in}

\begin{itemize}[leftmargin=*, itemsep=0.08in,
label={\color{brandblue}\textbullet}]

\item Very long sequential dependencies must be preserved.

\item Fine-grained control over memory is beneficial.

\item Large datasets are available.

\item Computational cost is less of a concern.

\item Slightly higher model capacity may improve performance.

\end{itemize}

\vspace{0.08in}

\textbf{Common Applications}

{\small
Financial time-series\\
Language modeling\\
Machine translation\\
Long sequence prediction
}

\end{minipage}

\vspace{0.2in}

\begin{infobox}

\textbf{Engineering Rule of Thumb}

Start with \textbf{GRU}. If the task involves very long-term
dependencies or GRU struggles to capture the required context,
consider switching to \textbf{LSTM}.

\end{infobox}

\vfill
\newpage

% ==========================================================
% SLIDE 9: THE BIG TAKEAWAY
% ==========================================================

\vspace*{0.35in}

\begin{center}

{\fontsize{14}{16}\selectfont\color{brandblue}\textbf{THE ENGINEERING TAKEAWAY}}

\vspace{0.3in}

\begin{minipage}{0.88\textwidth}
\centering

{\color{brandblue}\rule{6cm}{1pt}}

\vspace{0.2in}

{\fontsize{24}{28}\selectfont\bfseries\color{brandnavy}
GRU isn't designed to\\
outperform LSTM.
}

\vspace{0.15in}

% Now the explanation is set in a clean, readable size — no bold subheading
{\fontsize{12}{14}\selectfont\color{brandgray}
It is an elegant simplification.

\vspace{0.1in}
By merging states and reducing gates, GRU delivers competitive performance
with significantly less computational overhead.

\vspace{0.1in}
The real lesson: choosing an architecture isn't about finding the ``best'' model. It's about matching complexity perfectly to your constraints.
}

\vspace{0.2in}

{\color{brandblue}\rule{6cm}{1pt}}

\end{minipage}

\end{center}

\vfill

\begin{center}

\rule{0.6\textwidth}{0.5pt}

\vspace{0.12in}

{\fontsize{11}{13}\selectfont\color{brandgray}
Building AI systems by understanding the mathematics,\\
engineering, and intuition behind every architecture.
}

\vspace{0.12in}

{\fontsize{10}{12}\selectfont\color{brandgray}
\begin{tabular}{@{}l l}
GitHub: & \texttt{github.com/MrHeaven1y} \\
Portfolio: & \texttt{mrheaven1y.github.io} \\
LinkedIn: & \texttt{linkedin.com/in/dibayendu-mukherjee-bb897b267}
\end{tabular}
}

\vspace{0.12in}

{\fontsize{14}{16}\selectfont\color{brandnavy}\textbf{Dibayendu Mukherjee}}

\end{center}

\vspace*{0.2in}

\end{document}